\documentclass[reqno]{amsart}

\usepackage{amssymb}
\usepackage{amsthm}
\usepackage{amsmath}
\usepackage{enumitem}

\usepackage{tikz-cd}
\usepackage{leftindex}
\usepackage[capitalise]{cleveref}
\usepackage[all,2cell]{xy}
\UseAllTwocells
\usepackage{bm} %bold math symbols

\usepackage{marginnote}

\renewcommand{\descriptionlabel}[1]{\hspace\labelsep\upshape\bfseries
  #1}
\renewcommand{\labelitemi}{---}

\title{Groupoidal polygraphic homology} 

\author{Léonard Guetta}
\address{L\'eonard Guetta, Universiteit Utrecht, Utrecht Science Park,
Hans Freudenthalgebouw, ...}
\email{l.s.guetta@uu.nl}

\author{Fran\c{c}ois M\'etayer}
\address{Fran\c{c}ois M\'etayer, Universit\'e Paris Cit\'e, CNRS, IRIF, F-75013, Paris, France \&
Universit\'e Paris Nanterre}
\email{metayer@irif.fr}

\input{macro}


\begin{document}

\begin{abstract} 
We show that for a $1$-category $C$, the \ook k-polygraphic homology
of $C$ for any $k\geq 1$, that is taken with cofibrant resolutions in strict \ook
k-categories, does not depend on $k$ and is canonically isomorphic to
the homology of the classifying space of $C$. When $C$ is a groupoid,
we also show this for $k=0$. In particular, this means that the
classical homology of groups can be obtained by taking cofibrant
resolutions in strict \oo-groupoids.

In order to show these results, we develop the theory of discrete
Conduché fibrations in the category of strict \ook k-categories,
building on previous work by the first-named author.
\end{abstract}

\maketitle

\section{Introduction}\label{sec:intro}

The category $\oCat$ of strict, globular $\omega$-categories bears a
``canonical'' model structure structure introduced
in~\cite{lafontetal:folkms}. In this structure, the cofibrant objects
are of the form $\frcat P$ where $P$ is a polygraph
(see~\cite{abgmmm:polybk}) and $\frcat P$ denotes the \oo-category
freely generated by $P$. Let $\Chp$ be the category of chain complexes
of abelian groups in non-negative degree equipped with its usual
projective model structure. There is a left-Quillen abelianization
functor $\abel:\oCat\to \Chp$ (see~\ref{subsec:abel} below). Let then $\LL\abel:\Loc{\oCat}\to\Loc{\Chp}$ denote the
left-derived functor of $\abel$ between the corresponding homotopy
categories. The {\em polygraphic homology} of an \oo-category $C$
is by definition
\[\Hopol(C)=\LL\abel(C).\]
To compute this homology, we first produce
a cofibrant replacement of $C$, namely a polygraph $P$ together with an 
an acyclic fibration $p:\frcat P\to C$, then apply the functor $\abel$
directly to $\frcat P$, whence the name.

Now for each integer $k\geq 0$, an \ook k-category is an \oo-category
in which all cells are invertible in dimensions $>k$. The full
subcategory of $\oCat$ whose objects are \ook k-categories is denoted
by $\opCat k$. For instance $\opCat 0$ is just the category of strict
\oo-groupoids.
The canonical inclusion functor $\oemb k:\opCat k\to\oCat$ admits a left
adjoint $\opifun k:\oCat\to\opCat k$. It turns out that the canonical model 
structure on $\oCat$ transfers to $\opCat k$ via $\oemb k$ and that
$\opifun k$ is again left-Quillen. On the other hand, the functor
$\lambda$ easily factors through $\opifun k$. Thus we get an
abelianization functor $\abelk k:\opCat k\to\Chp$ such that
$\abel=\abelk k\circ\opifun k$, and $\LL\abel$ is naturally isomorphic
to $\LL\abelk k\circ \LL\opifun k$. This leads to define the
polygraphic homology of an \ook k-category $C$ by
\[\Hop{k}(C)=\LL\abelk k(C).\]

 
  
\input{basics} % sec:basics, sec:syntax
\input{conduche} % sec:basislift
\input{resolproj} % sec:projresol
\section{Polygraphic homology}
\begin{paragr}
  Recall that we have a so-called ``abelianization'' functor $\abel
  \colon \oCat \to \Chp$ from the category of $\oo$\nbd-categories to
  the category of chain complexes in non-negative degree, defined in
  the following way. For an $\oo$\nbd-category $C$ and $n\geq 0$
  \[
    \abel_n(C):=\Z C_n/\sim_n,
  \]
  where $\Z C_n$ is the free abelian group on the set of $n$\nbd-cells
  of $C$ and $\sim_n$ is the smallest equivalence relation compatible with
  the abelian group structure such that $x\comp{i} y
  \simeq _n x + y$ for any $x,y\in C_n$ that are
  $\comp{i}$-composables. The boundary operator is defined as
  \[
    \begin{aligned}
      \partial \colon \abel_{n+1}(C) &\to \abel_{n}(C)\\
      [x] &\mapsto [\target x] - [\source x].
    \end{aligned}
  \]
  One checks that this is well-defined, that we have $\partial
  \circ \partial=0$ and that the correspondence $C \mapsto
  \abel_{\bullet}(C)$ can be made functorial in an obvious
  way. Furthermore, $\abel$ admits a right adjoint. For
  details we refer to \cite{abgmmm:polybk}.

  By pre-composing with the inclusion functor, for any $k\geq 0$, we
  also get an abelinization functor
  \[
    \abelk{k} \colon \opCat{k} \overset{\emb{k}{\oo}}{\longrightarrow} \oCat \overset{\lambda}{\longrightarrow} \Chp.
  \]
  In particular, for $k=\oo$, we have $\abel=\abelk{\oo}$ by
  definition. Notice that since $\emb{k}{\oo}$ admits a right adjoint, so does $\abelk{k}$.
\end{paragr}

The abelianization of \ook k-polygraphs admits a useful explicit
description as stated in the following result.
\begin{proposition}
  Let $k \geq 0$, $P$ an \ook k-polygraph and $P_{n}$ a set of
  generators for each $n \geq 0$. Then $\abel_{\bullet}(\frgp P)$
  is canonically isomorphic to the chain complex
  \[
    \Z P_0 \overset{\partial}{\longleftarrow} \Z P_1
    \overset{\partial}{\longleftarrow} \dots,
  \]
  where $\partial$ is defined on generator by
  \[
    \partial[x] = [\target x] - [\source x].
  \]
\end{proposition}
\lgcomment{A-t-on besoin de cette description? Par ailleurs, la définition de l'opérateur bord dans cette proposition
  n'est pas très précise.}
\begin{proof}
\end{proof}
\begin{paragr}
  Let $0 \leq k
  \leq k' \leq \omega$. We are now going to analyze the behaviour of
  the abelianization functor with respect to the left adjoint $\pifun{k'}{k} \colon
  \opCat{k'} \to \opCat{k}$ of the inclusion functor. By definition,
  we have a commutative triangle
  \[
    \begin{tikzcd}[column sep=small]
      \opCat{k}\ar[rr,"\emb{k}{k'}"]\ar[rd,"\abelk{k}"']&& \opCat{k'}\ar[dl,"\abelk{k'}"] \\
      &\Chp&.
  \end{tikzcd}
  \]
  Now, consider the
  following, \emph{a priori} not commutative, triangle
  \[
    \begin{tikzcd}[column sep=small]
      \opCat{k'}\ar[rr,"\pifun{k'}{k}"]\ar[rd,"\abelk{k'}"']&& \opCat{k}\ar[dl,"\abelk{k}"] \\
    &\Chp&.
  \end{tikzcd}
\]
The unit defines a natural transformation
\[
 \abelk{k'}\Longrightarrow \abelk{k'} \emb{k}{k'}\pifun{k'}{k}= \abelk{k}\pifun{k'}{k}
\]
\end{paragr}
\begin{lemma}\label{lemma:abelpi}
  The above natural transformation $\abelk{k'} \Rightarrow \abelk{k}\pifun{k'}{k}
  $ is an isomorphism. 
\end{lemma}
\begin{proof}
\end{proof}
\begin{paragr}
  Recall that the category $\oCat$ admits a monoidal category
  structure with tensor product the so-called ``Gray tensor product'',
  simply denoted by $\gray$, and with unit $\Dsk{0}$ (see for example
  \cite{}). More generally, for any $0 \leq k \leq \omega$, we have
  the induced tensor product $\grayk{k}$
\end{paragr}
\begin{proposition}
  For any $0 \leq k \leq \omega$, the abelianization functor
  \[
    \abelk{k} \colon \opCat{k} \to \Chp
  \]
  is strong monoidal with respect to the Gray tensor product $\grayk{k}$ on
  $\opCat{k}$ and the usual tensor product of chain complexes on $\Chp$.
\end{proposition}
\begin{proof}
  The case $k=\omega$ is standard (see for example \cite{}). The
  general case follows from  $\abelk{k}\simeq
  \abelk{\omega}\pifun{\oo}{k}$ (\cref{lemma:abelpi}) and the fact
  that $\abelk{\omega}$ and $\pifun{\oo}{k}$ are strong monoidal.
\end{proof}
\begin{proposition}
  For any $k \geq 0$, the abelianization functor
  \[
    \abelk{k} \colon \opCat{k} \to \Chp
  \]
  is left Quillen, when $\npCat{\oo}{k}$ is equipped with the folk
  model structure and $\Chp$ with the projective model structure.
\end{proposition}
\begin{proof}
    The case $k=\omega$ is \cite[Theorem 22.2.7]{abgmmm:polybk}. For the
  general case, since the folk model structure on $\opCat{k}$ is right induced from
  the inclusion functor $\emb{k}{\omega} \colon \opCat{k} \to \oCat
  $, this means that a set of generating cofibrations (resp.\ trivial
  cofibrations) is given by the image by $\pifun{k}{\omega} \colon \oCat
  \to \opCat{k}$ of generating cofibrations of
  $\pifun{k}{\omega}$. The result follows then immediatly from \cref{lemma:abelpi}.
\end{proof}
\begin{paragr}
  The previous proposition contains an important subtlety. % For any $0
  % \leq k \leq k' \leq \omega$, the inclusion functor
  % $\embed{k}{k'}\colon \opCat{k} \to \opCat{k}$ is
  Even though for any  $0 \leq k \leq k' \leq \omega$ we have a
  commutative triangle,
  \[
    \begin{tikzcd}[column sep=small]
      \opCat{k}\ar[rr,"\emb{k}{k'}"]\ar[rd,"\abelk{k}"']&& \opCat{k'}\ar[dl,"\abelk{k'}"] \\
      &\Chp&,
  \end{tikzcd}
\]
and even though $\emb{k}{k'}$ passes to the homotopy categories, as
it trivially preserves the weak equivalences, this doesn't \emph{a
  priori} imply that the following triangle at the level of homotopy categories
\[
  \begin{tikzcd}[column sep=small]
    \Loc{\opCat{k}}\ar[rr,"\emb{k}{k'}"]\ar[rd,"\LL \abelk{k}"']&&
    \Loc{\opCat{k'}}\ar[dl,"\LL \abelk{k'}"] \\
    &\Loc{\Chp}&,
  \end{tikzcd}
\]
is also commutative, and the universal property of left derived functors
only yields a natural transformation
\begin{equation}\label{transnat}
  \LL \abelk{k'}\circ \emb{k}{k'} \Rightarrow \LL \abelk{k},
\end{equation}
which is \emph{a priori} not an isomorphism (and we do not know if it
is in general, see \cref{openquestion}.)

More concretly, this natural transformation can be defined as follows:
let $C$ be an \ook k\nbd-category and let
\[
  P\to C
\]
be a cofibrant resolution in $\opCat{k}$. Since $\emb{k}{k'}$ does
not preserve cofibrant objects in general, we need to take a cofibrant
resolution in $\opCat{k'}$
\[
  P' \to \emb{k}{k'}P.
\]
The above natural transformation is then
\[
  \LL \abelk{k'}(\emb{k}{k'}(C))=\abelk{k'}(P') \longrightarrow
  \abelk{k'}(\emb{k}{k'}(P))=\abelk{k}(P)=\LL \abelk{k}(C).
\]
\end{paragr}

\begin{definition}
  Let $k\geq 0$. The \emph{$\pair \oo k$\nbd-polygraphic homology} of
  an \ook k\nbd-category $C$ is defined as
  \[
    \Hop{k}_{\bullet}(C):=\Ho_{\bullet}(\LL \abelk{k} (C)).
  \]
\end{definition}

\begin{paragr}
  For $0 \leq k \leq k'\leq \omega$ and an \ook k\nbd-category $C$, we shall
  simply write $\Hop{k'}_{\bullet}(C)$ instead of
  $\Hop{k'}_{\bullet}(\emb{k}{k'}(C))$.
  In particular, the natural transformation \eqref{transnat} becomes
  \[
    \Hop{k'}_{\bullet}(C) \to \Hop{k}_{\bullet}(C).
  \]
\end{paragr}

\begin{openquestion}\label{openquestion} With the same notation as
  above, is the canonical morphism
  \[
     \Hop{k'}_{\bullet}(C) \to \Hop{k}_{\bullet}(C)
   \]
   an isomorphism of $\Loc{\Chp}$?
\end{openquestion}
As we shall see at the end of this section, we answer positively to
the above open question in a very particular subcase. 

\begin{theorem}
  Let $C$ be a $\pk 1 i$\nbd-category with $i \in \{0,1\}$. Then, for any $i
  \leq k \leq \oo$ there is an isomorphism
  \[
    \Hop{k}_{\bullet}(C) \to \Ho_{\bullet}(C)
  \]
  natural in $C$. Furthermore, this isomorphism can be chosen
  naturally in $k$, meaning that for $i \leq k \leq k' \leq \oo$, the
  triangle
  \[
    \begin{tikzcd}[column sep=small]
         \Hop{k'}_{\bullet}(C)\ar[dr]
         \ar[rr]&&\Hop{k}_{\bullet}(C)\ar[dl]\\
         &\Ho_{\bullet}(C)&
    \end{tikzcd}
  \]
  is commutative. 
\end{theorem}
\begin{proof}
  We begin by the first part of the statement. Because $k \geq i$, we
  can see $C$ as an $\ook k$\nbd-category, via the canonical
  inclusion $\npCat{1}{i} \to \opCat{k}$, which we omit to denote. Let $P \to C$ be a
  polygraphic resolution in $\opCat{k}$. By 
\end{proof}
%%% Local Variables:
%%% mode: LaTeX
%%% TeX-master: "main"
%%% End:
 % sec:polhom
\section{Recollection on homotopy kan extensions in Model categories}
\subsection{Projective model structure}\label{projms}

\begin{paragr}\label{paragr:htpykan} Let $\M$ be a cofibrantly generated model
category. Recall that for any small category $I$, we can equip the
category of functors $\fun{I}{\M}$ with the \emph{projective model
structure}, where the weak equivalences and the fibrations are these
natural transformations which are pointwise weak equivalences and
fibrations of $\M$ respectively. Any functor $u \colon I \to J$
between small categories induces by precomposition a functor
  \begin{align*} u^* \colon \fun{J}{\M} &\to \fun{I}{\M}\\ X &\mapsto
X\circ u.
  \end{align*} This functor admits a left adjoint $u_!
\colon\fun{I}{\M} \to \fun{J}{\M} $ given by the pointwise Kan
extension formula:
  \[ u_!(X)(j)=\ilim_{\tr{I}{j}}X\vert_{\tr{I}{j}}.
  \] By construction, $u^*$ preserves the weak equivalences and the
fibrations of the projective model structure. Hence, the adjunction
$u_! \dashv u^*$ is a Quillen adjunction and we get an adjunction at
the level of homotopy categories
  \[ \LL u_! \colon \Loc{\fun{I}{\M}} \leftrightarrows
\Loc{\fun{J}{\M}} \colon u^*.
  \] Here, we kept the same notation $u^*$ at the level of homotopy
categories (instead of, say, $\mathbb{R} u^*$) because the functor
$u^*$ preserves weak equivalences and thus we have a commutative
square
  \[
    \begin{tikzcd} \fun{J}{\M} \ar[d] \ar[r,"u^*"] &
\fun{I}{\M}\ar[d]\\ \Loc{\fun{J}{\M}} \ar[r,"u^*"] &
\Loc{\fun{I}{\M}},
    \end{tikzcd}
  \] where the vertical arrows are the localization functors.
\end{paragr}
\begin{paragr} We denote by $\Term$ the terminal category and make the
canonical identification $\fun{\Term}{\M}=\M$. For a small category
$I$, we denote by $p_I \colon I \to \Term$ the unique such
functor. The functor $p_I^* \colon \M \to \fun{I}{\M}$ is the
``diagonal'' functor that associates to an object $X$ of $\M$, the
constant diagram $I \to \M$ with value $X$. Hence, $p_{I!}$ is nothing
but the colimit functor $\ilim_I \colon \fun{I}{\M} \to \M$, and we
shall use the two notations interchangeably. Consequently, $\LL
p_{I!}$ is nothing but the \emph{homotopy colimit functor}
  \[ \hoilim_I \colon \Loc{\fun{I}{\M}} \to \Loc{\M}.
  \]
\end{paragr}
\begin{paragr} Let $I$ be a small category and $i \in \Ob(I)$ an
object of $I$. This can be seen as a functor $i \colon \Term \to I$,
which induces a ``evaluation at $i$'' functor
  \begin{align*} i^* \colon \fun{I}{\M} &\to \M\\ X &\mapsto X(i).
  \end{align*} It is a standard fact of category theory that the
family of functors
  \[\left(i^*\colon \fun{I}{\M} \to \M\right)_{i\in \Ob(I)}\] is
\emph{conservative}, meaning that a morphism $\alpha \colon X \to Y$
of $\fun{I}{\M}$ is an isomorphism if and only if for each object $i$
of $I$, $i^*(\alpha)=\alpha_i$ is an isomorphism. The same is true at
the level of homotopy categories as we recall in the following result.
\end{paragr}
\begin{proposition} Let $\M$ be a (cofibrantly generated) model
category and $I$ a small category. The family of functors
  \[ \left(i^* \colon \Loc{\fun{I}{\M}}\to \Loc{\M}\right)_{i \in
\Ob(I)}
  \] is conservative.
\end{proposition}
\begin{proof}
  
\end{proof}
\subsection{Homotopy cocontinuous functors}\label{hptycocts}
\begin{paragr} Let $F \colon \M \to \M'$ be a left Quillen functor
between (cofibrantly generated) model categories. For any small
category $I$, we consider the functor induced by post-composition
  \begin{align*} F^I \colon \fun{I}{\M} &\to \fun{I}{\M'}\\ X &\mapsto
F\circ X.
  \end{align*} This functor is also a left Quillen functor with
respect to the projective model structures. Hence, it induces a left
derived functor
  \[ \LL F^I \colon \Loc{\fun{I}{\M}}\to \Loc{\fun{I}{\M'}}.
  \]

  Now, let $u \colon I \to J$ be a functor between small
categories. By construction, we have a commutative square
  \[
    \begin{tikzcd} \fun{J}{\M} \ar[r,"F^J"] \ar[d,"u^*"']&
\fun{J}{\M'} \ar[d,"u^*"]\\ \fun{I}{\M} \ar[r,"F^I"'] & \fun{I}{\M'},
    \end{tikzcd}
  \] and, because $F$ is a left adjoint, the following square is
commutative up to a canonical isomorphism
  \[
    \begin{tikzcd} \fun{I}{\M} \ar[r,"F^I"] \ar[d,"u_!"']&
\fun{I}{\M'} \ar[d,"u_!"]\\ \fun{J}{\M} \ar[r,"F^J"'] & \fun{J}{\M'}.
\ar[from=1-1,to=2-2,phantom,"\simeq" description]
    \end{tikzcd}
  \]
  This fact is usually referred to as ``$F$ preserves Kan
extension''. We recall now the following result which says that the
same holds at the level of homotopy categories for a left Quillen
functor.
\end{paragr}
\begin{proposition} Let $\M$ and $\M'$ be two (cofibrantly generated)
  model categories, $F \colon \M \to \M'$ be a left Quillen functor and
  $u \colon I \to J$ a functor between small categories. Then the
  squares
  \[
    \begin{tikzcd} \Loc{\fun{J}{\M}} \ar[r,"\LL F^J"] \ar[d,"u^*"']&
      \Loc{\fun{J}{\M'}} \ar[d,"u^*"]\\ \Loc{\fun{I}{\M}} \ar[r,"\LL F^I"']
      & \Loc{\fun{I}{\M'}},
    \end{tikzcd}
  \] and
  \[
    \begin{tikzcd} \Loc{ \fun{I}{\M}} \ar[r,"\LL F^I"] \ar[d,"\LL
      u_!"']& \Loc{\fun{I}{\M'}} \ar[d,"\LL u_!"]\\ \Loc{\fun{J}{\M}}
      \ar[r,"\LL F^J"'] & \Loc{\fun{J}{\M'}},
    \end{tikzcd}
  \] are commutative up to a canonical isomorphism.
\end{proposition}
\begin{proof}
\end{proof}
\subsection{Abstract homology of small categories}\label{abstracthom}
\begin{definition}\label{def:abstracthom}
  Let $\M$ be a (cofibrantly generated) model category. For a small
  category $I$ and $X$ an object of $\fun{I}{\M}$, we define the \emph{$\M$\nbd-valued
    homology of $I$ with coefficient $X$} as the object of $\Loc{\M}$
  \[
    \Ho_{\M}(I,X):=\LL p_{I!}X=\hoilim_IX.
  \]
\end{definition}
\begin{remark}
  Note that the use of the word ``homology'' here is arguably slightly
  abusive as we don't make any hypothesis of additivity or stability
  property of any kind on (the localization of) $\M$. 
\end{remark}
\begin{paragr}
 Obviously, for a fixed object small category $I$, the previous
 construction yields a functor
 \[
   \Ho_{\M}(I,-) \colon \Loc{\fun{I}{\M}}\to \Loc{\M}.
 \]
 Let us now explore the functoriality of $\Ho_{\M}(I,X)$ in $I$. Let
  $u \colon I \to J$ be a functor between small categories. We have a
  commutative triangle
  \[
    \begin{tikzcd}[column sep=small]
      I \ar[rr,"u"]\ar[rd,"p_I"'] && J\ar[ld,"p_J"] \\
      &\Term&
    \end{tikzcd}
  \]
  which induces a commutative triangle
  \[
    \begin{tikzcd}[column sep=small]
      \Loc{\fun{I}{\M}}  && \Loc{\fun{J}{\M}}\ar[ll,"u^*"'] \\
      &\Loc{\M}.\ar[lu,"p_I^*"]\ar[ru,"p_J^*"']&
    \end{tikzcd}
  \]
  Hence, by uniqueness of adjoints, this yields a triangle commutative
  up to a canonical isomorphism
  \[
        \begin{tikzcd}[column sep=small]
      \Loc{\fun{I}{\M}} \ar[rr,"\LL u_!"]\ar[rd,"\LL p_{I!}"'] &&
      \Loc{\fun{J}{\M}}\ar[ld,"\LL p_{J!}"] \\
      &\Loc{\M}.&
    \end{tikzcd}
  \]
  Now let $X$ be an object of $\fun{J}{\M}$. Then we have a morphism
  \[
    \LL p_{I!}u^*X \simeq \LL p_{J!} \LL u_! u^*X\to \LL p_{J!}X,
  \]
  where the morphism of the right is the co-unit of the ajdunction
  $\LL u_! \dashv u^*$. On other words, this defines a morphism
  \[
  \Ho_{\M}(u,X)\colon  \Ho_{\M}(I,u^*X)\to \Ho_{\M}(J,X).
  \]
  We then have the following functoriality result.
\end{paragr}
\begin{proposition}\label{prop:functorialityhomology}
  Let $\M$ a (cofibrantly generated) model category, $K$ a small
  category and $X$ an object of $\fun{K}{\M}$. We have a functor
  \[
      \Ho_{\M}(-,X) \colon \tr{\Cat}{K} \to \Loc{\M},
    \]
    whose action on object is 
  \[
    (I,p \colon I \to K) \longmapsto \Ho_{\M}(I,p^*X)
  \]
  and whose action on morphisms is
  \[
    u\colon (I,p)\to (J,q) \longmapsto \Ho_{\M}(u,q^*X)\colon \Ho_{\M}(I,p^*X) \to \Ho_{\M}(J,q^*X).
\]
 
\end{proposition}
\begin{proof}
\end{proof}

\begin{paragr}
  The main case of interest of the previous proposition will be in the
  case that $K=\Term$ is the terminal category. Then $F$ is simply an
  object $X$ of $\M$ and we can make the identification
  $\tr{\Cat}{K}\simeq \Cat$. We speak of the (abstract) homology of a
  small category $I$ with \emph{constant} coefficient $X$ for
  \[
    \Ho(I,p_I^*X).
  \]
  \cref{prop:functorialityhomology} in this case yields the following corollary.
\end{paragr}
\begin{corollary}
  Let $\M$ be a (cofibrantly generated) model category and $X$ an
  object of $\M$. We have a functor
  \[
    \Ho(-,X) \colon \Cat \to \Loc{\M},
  \]
  whose action on objects is
  \[
    I \longmapsto \Ho_{\M}(I,p_I^*X),
  \]
  and whose action on morphisms is
  \[
    u \colon I \to J \longmapsto \Ho_{\M}(u,X) \colon
    \Ho_{\M}(I,p_I^*X) \to \Ho_{\M}(J,p_J^*X).
  \]
\end{corollary}
\subsection{Homology of categories}
\begin{paragr}\label{paragr:homalg}
  For $C$ a small category, we denote by $\LMod{C}$ the category of
  \emph{right $C$\nbd-modules}, that is of functors
  \[
    C \to \Ab
  \]
  from $C$ to the category of abelian groups, and of natural
  transformations between them. By definition, the category of \emph{right}
  $C$\nbd-modules is the category
  \[
    \RMod{C}:=\LMod{C^{\op}}
  \]
  and everything that follows is dualized accordingly. 
  
  We denote by $\Chp(C)$ the category of chain complexes
  of left $C$\nbd-modules in non-negative degree. For $C=\Term$ the
  terminal category, we use the notation $\Chp(\Z)$ instead. 
  We denote by $\Derp{C}$ the \emph{derived
 category of $C$} in non-negative degree, that is the
localization of $\Chp(C)$ with respect to
 quasi-isomorphisms (and we use the notation $\Derp{\Z}$ when
 $C=\Term$ is the terminal category). This localization is controlled by the
  so-called \emph{projective} model structure on $\Chp(C)$, which is
  cofibrantly generated,
  characterized as follows:
  \begin{itemize}
  \item the weak equivalences are the quasi-isomorphisms,
  \item the fibrations are the levelwise epimorphisms in each positive
    degree,
  \item the cofibrations are the levelwise monomorphisms with
    projective cokernel in each non-negative degree.
  \end{itemize}
  In particular, cofibrant objects are the chain complexes (in
  non-negative degree) of projective
  $C$\nbd-modules. The following lemma is useful to
  construct and recognize such objects.
\end{paragr}
\begin{lemma}\label{lemma:freeisproj}
  Let $c$ be an object of $C$. The left $C$\nbd-module $\Z[\homset{C}{c}{-}]$
  \[
    \begin{aligned}
      C &\to \Ab\\
      d &\mapsto \Z[\homset{C}{c}{d}],
    \end{aligned}
  \]
  is a projective object of $\LMod{C}$.
\end{lemma}
\begin{proof}
  This follows straightforwardly from the Yoneda lemma.
\end{proof}

\begin{paragr}
   Notice that for a small category $C$, we have a canonical
   \emph{isomorphism} of categories
  \[
    \Chp(C)\simeq \fun{C}{\Chp(\Z)},
  \]
  and we shall make this identification implictly. Note that this identification is compatible with
  model structures in the sense that the projective model structure
  on $\fun{C}{\Chp(\Z)}$ in the sense of \ref{paragr:htpykan} is the
  same, via the previous isomorphism, as the projective model
  structure on $\Chp(C)$ in the sense of \ref{paragr:homalg}.

  This allows us to use all the results from the sections
  \ref{projms}, \ref{hptycocts} and \ref{abstracthom}. 
\end{paragr}
\begin{definition}
  Let $C$ be a small category and $X$ a left $C$\nbd-module. We define
  the \emph{homology of $C$ with coefficient $X$}, and denote by
  \[
    \Ho(C,X),
  \]
  the $\Chp(\Z)$\nbd-valued homology of $C$ with coefficient $X$
  (\cref{def:abstracthom}).
\end{definition}
%%% Local Variables:
%%% mode: LaTeX
%%% TeX-master: "main"
%%% End:
 % sec:htpykan
\section{Comparison theorem}\label{sec:comparison}

\subsection{The two fundamental properties}
\begin{paragr}\label{paragr:colimtr}
  Let $C$ be a $1$\nbd-category. Recall from \ref{paragr:pullback_resol} that, given an $\ook
  k$\nbd-category
  $X$ and an $\ook k$\nbd-functor $p \colon X \to C$, with $1 \leq k \leq \oo$, we defined a functor
    \[
    \begin{aligned}
      \tr{X}{-} \colon C &\to \opCat{k} \\
      c &\mapsto \tr{X}{c}.
    \end{aligned}
  \]
    This
    construction is functorial in $p\colon X \to C$ as it canonically defines a functor
  \begin{align*}
    \tr{\opCat{k}}{C} &\to \funopCat{k}{C}\\
    (X, X \overset{p}{\to} C) &\mapsto \left(c \mapsto \tr{X}{c}\right).
  \end{align*}
  Notice that for $(X, X \overset{p}{\to} C)$ in $\tr{\opCat{k}}{C}$,
  we have a cocone
  \[
    (\tr{X}{c} \to X)_{c \in \Ob{C}},
  \]
  hence
  a canonical morphism
  \[
    \ilim_{c \in C}\tr{X}{c} \to X,
  \]
  which is natural in the sense that it defines a natural
  transformation from the functor
  \begin{align*}
    \tr{\opCat{k}}{C} &\to \opCat{k}\\
    (X, X \overset{p}{\to} C) &\mapsto \ilim_{c \in C} \tr{X}{c},
  \end{align*}
  to the projection functor
    \begin{align*}
    \tr{\opCat{k}}{C} &\to \opCat{k}\\
    (X, X \overset{p}{\to} C) &\mapsto X.
    \end{align*}
    Finally, if we take $C$ a groupoid and $k=0$, then everything said
    above still makes sense.
\end{paragr}
\begin{lemma}\label{lemma:colim}
  The natural morphism of $\opCat{k}$
  \[
     \ilim_{c \in C}\tr{X}{c} \to X,
   \]
   is an isomorphism.
\end{lemma}
\begin{proof}
The case $k=\oo$ is \cite[Lemma 7.5]{guetta:homcat}. The general case
follows from the fact that $\opCat{k} \to \oCat$
  preserves colimits.
\end{proof}
\begin{proposition}\label{prop:hocolimtr}
  Let $C$ be a $1$\nbd-category. Then for any $1 \leq k \leq \oo$, the canonical morphism of
  $\Loc{\opCat{k}}$
  \[
    \hocolim_{c\in C }\tr{C}{c} \to C
    \]
  is an isomorphism. If $C$ is a groupoid, then this also works for
  $k=0$.
\end{proposition}
\begin{proof}
  Let $\frgp{P} \to C$ be a polygraphic resolution of $C$ in
  $\opCat{k}$ and \[\tr{\frgp P\!}{-} \colon C \to \opCat{k}\] obtained
  as in section \ref{subsec:polres}. We get a commutative diagram in
  $\Loc{\opCat{k}}$
  \[
    \begin{tikzcd}
      \displaystyle\hocolim_{c \in C}\tr{\frgp P\!}{c} \ar[d] \ar[r]& \displaystyle\ilim_{c \in
        C}\tr{\frgp P\!}{c} \ar[d] \ar[r] & \frgp P \ar[d] \\
      \displaystyle\hocolim_{c \in C}\tr{C}{c} \ar[r]&\displaystyle\ilim_{c \in C}\tr{C}{c}\ar[r] & C.
    \end{tikzcd}
  \]
  Note that both $\ilim_{c \in C}$ in the middle are the images by the
  localisation functor \[\opCat{k} \to \Loc{\opCat{k}}\]
  of the actual
  colimits taken in the category $\opCat{k}$. The right vertical
  arrow is an isomorphism by hypothesis, the horizontal arrows of the
  square on the right are isomorphisms by \cref{lemma:colim}. By
  stability of trivial fibrations under pullback, the natural transformation
  $\tr{\frgp P\!}{-} \to \tr{C}{-}$ is a pointwise weak equivalence. Since homotopy colimits
  send pointwise weak equivalences to isomorphisms (in the homotopy
  category), this proves that the left vertical arrow is an
  isomorphism. Finally, the top left
  horizontal arrow is an isomorphism
  because $\tr{\frgp P\!}{-}$ is cofibrant by \cref{thm:projresol} (or
  \cref{thm:projresol_grp} in the case $k=0$). By applying
  two-out-of-three twice, we get that all the remaining arrows of this diagram
  are isomorphisms too. 
\end{proof}
\begin{lemma}\label{lemma:liftoplax}
  Let $0 \leq k \leq \omega$ and consider a diagram as below in $\opCat{k}$, 
  \[
    \begin{tikzcd}[row sep=large]
      A' \ar[d,"p"]\ar[r,bend left,"u'"] \ar[r,bend right,"v'"']& B' \ar[d,"q"] \\
      A \ar[r,bend left,"u",""{name=A,below}] \ar[r,bend
      right,"v"',""{name=B,above}]& B,
      \ar[from=A,to=B,Rightarrow,"\alpha"]
    \end{tikzcd}
  \]
  where the front and back faces are commutative. If $q$ is a trivial
  fibration and $A'$ is cofibrant in the folk model structure on
  $\opCat{k}$, then there exists an oplax transformation $\alpha'
  \colon u' \Rightarrow v'$ making the whole diagram commute, which means
  that
  \[
    q\wsk \alpha' = \alpha \wsk p.
  \]
\end{lemma}
\begin{proof}
  The data of the given diagram amounts to the following commutative
  diagram in $\opCat{k}$
  \[
    \begin{tikzcd}
      \Sph{0} \grayk{k} A' \ar[r]\ar[d] \ar[rr,bend
      left,"{(u',v')}"]& \opifun{k}(\Dsk{1}) \grayk{k} A'
      \ar[d]& B' \ar[d,"q"] \\
      \Sph{0} \grayk{k} A \ar[r] & \opifun{k}(\Dsk{1}) \grayk{k} A \ar[r,"\alpha"]
      & B.
    \end{tikzcd}
  \]
  Since $A'$ is cofibrant and $\grayk{k}$ makes the folk model
  structure on $\opCat{k}$ a monoidal model structure \cite{ara_lucas:folmon}, the
  arrow $ \Sph{0}  \grayk{k} A' \to  \opifun{k}(\Dsk{1}) \grayk{k} A'$ is a
  cofibration. Since $q$ is a trivial fibration, there exists a
  morphism $\alpha' \colon \opifun{k}(\Dsk{1})  \grayk{k} A' \to B'$ making the
  whole diagram commute.
\end{proof}
\begin{proposition}\label{prop:hlgycontrcat}
  Let $C$ be a small category with a terminal object. Then, for any $1
  \leq k \leq \oo$ the
  canonical morphism of $\Derp{\Z}$
  \[
    \Hop{k}(C) \to \Z,
  \]
  induced by the canonical functor $C \to \Term$, is an
  isomorphism. If $C$ is a groupoid, then it is also an isomorphism
  for $k=0$.
\end{proposition}
\begin{proof}
  Let us denote by $c_0$ the terminal object of $C$. We then have a
  natural transformation
    \[
    \begin{tikzcd}
      C \ar[r,bend left,"\Unit{C}",""{name=A,below}]
      \ar[r,bend right,"\cst{c_0}"',""{name=B,above}] & C
      \ar[from=A,to=B,Rightarrow]
    \end{tikzcd}
  \]
  from the identity functor on $C$ to the constant functor with value
  $c_0$.  By \ref{paragr:gray}, this defines an oplax
  transformation. Now, let
  \[
    \frgp{P} \to C
  \]
  be a polygraphic resolution of $C$. Since trivial fibrations are
  surjective on objects, we may choose an object $x_0$ of $P$ whose image in
  $C$ is $c_0$. Thus, we have a diagram
  \[
    \begin{tikzcd}[row sep=huge]
      \frgp{P} \ar[d]\ar[r,bend left,"\Unit{\frgp{P}}"] \ar[r,bend
      right,"\cst{x_0}"']& \frgp{P} \ar[d] \\
      C \ar[r,bend left,"\Unit{C}",""{name=A,below}] \ar[r,bend
      right,"\cst{c_0}"',""{name=B,above}]& C,
      \ar[from=A,to=B,Rightarrow]
    \end{tikzcd}
  \]
  such that the front and back faces commute. Since
  $\frgp{P}$ is cofibrant, and since the vertical arrows are trivial
  fibrations, we can apply \cref{lemma:liftoplax}, which yields an
  oplax transformation from the identity \oo\nbd-functor on $P$ to
  the constant \oo\nbd-functor with value $x_0$.

  Finally, applying \cref{cor:abeltrans}, we get a chain homotopy
  between the identity morphism \[\abelk{k}(P) \to
    \abelk{k}(P)\] and the composition
  \[\abelk{k}(P) \to \Z \to
  \abelk{k}(P),\] which proves that  $\abelk{k}(P) \to
  \Z$ is a quasi-isomorphism.
\end{proof}
\subsection{The comparison theorem}
\begin{theorem}\label{thm:polhom}
  Let $C$ be a $1$\nbd-category. Then, for any $1
  \leq k \leq \oo$, there is a canonical isomorphism
  \[
    \Hop{k}(C) \to \Ho(C)
  \]
  natural in $C$. In particular, for $k=\oo$, we get an isomorphism
  $\Hopol(C) \to \Ho(C)$. Furthermore, for any $k\geq 1$, these
  isomorphisms make the following triangle commute
  \[
    \begin{tikzcd}[column sep=small]      
         \Hopol(C)\ar[dr]
         \ar[rr,"\homnat_C"]&&\Hop{k}(C)\ar[dl]\\
         &\Ho(C).&
    \end{tikzcd}
  \]
\end{theorem}

\begin{proof}
  Let us start with the first part of the theorem. Recall that if $F
  \colon \M \to \M'$ is a functor, and $C$ is a small category we
  denote by $\fun{C}{F} \colon \fun{C}{\M} \to \fun{C}{\M'}$ the
  functor induced by postcomposition. Consider the
  following zig-zag in $\Derp{\Z}$
  \[
    \begin{tikzcd}
      \hocolim_{C}p^*_C(\LL \abelk{k}(\Term))&&\\
    \hocolim_{C}\LL \fun{C}{\abelk{k}}(p_C^*\Term) \ar[u,"(a)"]&\ar[l,"(b)"'] \hocolim_{C}\LL\fun{C}{
     \abelk{k}}(\tr{C}{-})\ar[r,"(c)"]&  \LL
   \abelk{k}(\hocolim_{C}\tr{C}{-})\ar[d,"(d)"]\\
   && \LL\abelk{k}(C),
  \end{tikzcd}
  \]
  where
  \begin{itemize}
    \item $(a)$ and $(c)$ are the universal morphisms constructed in \ref{paragr:univmor},
    \item $(b)$ comes from the fact that $p_C^*\Term$ is the terminal
      object of $\funopCat{k}{C}$,
      \item $(d)$ is induced by the cocone $(\tr{C}{c}\to
  C)_{c \in \Ob(C)}$ (see \ref{paragr:colimtr}).
\end{itemize}
All these morphisms of $\Derp{\Z}$ are in fact isomorphisms:
\begin{itemize}
  \item $(a)$ by \eqref{eq:morphder} of \cref{prop:htpycocts} and the fact that $\abelk{k}$
    is a left Quillen functor (\cref{prop:abelquillen}),
      \item $(b)$ by \cref{prop:der2}, \cref{prop:hlgycontrcat} and the
    fact that for each object $c$ of $C$ the category $\tr{C}{c}$ has
    a terminal object,
  \item $(c)$ by \eqref{eq:cocontinuous} of \cref{prop:htpycocts} and the fact that $\abelk{k}$
    is a left Quillen functor,
  \item $(d)$ by \cref{prop:hocolimtr}.
  \end{itemize}
  Finally, since by
  \cref{paragr:abeldisk} and the fact that $\Dsk{0}=\Term$, we have
  \[
    \LL \abelk{k}(\Term) \cong \Z.
  \]
  Since, by definition, we have
  \[
    \Ho(C)=\hocolim_{C}p^*_C(\Z),
    \] this yields an isomorphism
  \[
    \Ho(C) \cong \Hop{k}(C),
  \]
  which is natural in $C$ by construction, as all the morphisms in the
  previous zig-zag are natural in $C$.

  For the second part, consider the following diagram in $\Derp{\Z}$
  \[
    \begin{tikzcd}
      \LL\abel(C)\ar[r] & \LL\abelk{k}(C)\\
       \LL\abel(\hocolim_{C}\tr{C}{-})\ar[r]\ar[u]&
       \LL\abelk{k}(\hocolim_{C}\tr{C}{-})\ar[u]\\
       \hocolim_{C}\fun{C}{\LL
     \abel}(\tr{C}{-})\ar[r] \ar[u]\ar[d]& \hocolim_{C}\fun{C}{\LL
     \abelk{k}}(\tr{C}{-})\ar[u]\ar[d]\\
   \hocolim_{C}\LL \fun{C}{\abel}(p_C^*\Term)\ar[r] \ar[d]& \hocolim_{C}\LL
   \fun{C}{\abelk{k}}(p_C^*\Term)\ar[d] \\
    \hocolim_{C}p^*_C(\LL \abel(\Term))\ar[r]& \hocolim_{C}p^*_C(\LL
    \abelk{k}(\Term)),
    \ar[from=1-1,to=2-2,phantom,description, "(A)"]
    \ar[from=2-1,to=3-2,phantom,description,"(B)"]
    \ar[from=3-1,to=4-2,phantom,description,"(C)"]
    \ar[from=4-1,to=5-2,phantom,description,"(D)"]
  \end{tikzcd}
\]
where the horizontal arrows are induced by the natural transformation
\eqref{transnat} from \cref{paragr:subtlety}
\[
    \homnat \colon \LL \abel\circ \oemb{k} \Rightarrow \LL \abelk{k}.
  \]
  Note that for simplicity we have not made the functor
  $\oemb{k} \colon \Loc{\opCat{k}}\to \Loc{\oCat}$ appear in the above
  diagram. This naturality implies that the squares $(A)$ and $(C)$ are
  commutative. Furthermore, the commutativity of the squares $(B)$ and
  $(D)$ follows from the naturality in $F$ of \eqref{eq:morphder} and
  \eqref{eq:cocontinuous} from \cref{prop:htpycocts}.

  To conclude, we only need to notice that the following triangle is
  commutative as an immediate computation shows
  \[
    \begin{tikzcd}
      \LL \abel(\Term) \ar[d,"\cong"']\ar[r] & \LL \abelk{k}(\Term)\ar[dl,"\cong"]\\
      \Z.&
    \end{tikzcd}
  \]
\end{proof}
\begin{corollary}
  For any $1$\nbd-category $C$ and any $1 \leq k \leq \oo$, the
  canonical natural morphism
  \[
   \homnat_C \colon \Hopol(C) \to \Hop{k}(C)
  \]
  is an isomorphism. If $C$ is a groupoid, this also holds for $k=0$.
\end{corollary}
\begin{remark}
  Adapting \emph{mutatis mutandis} the proof of
  \cref{thm:polhom}, one can prove more generally that, for a
  $1$-category $C$ and for any $1 \leq k \leq k' \leq \oo$, there is a canonical natural isomorphism
  \[
     \Hop{k'\!}(C)\cong  \Hop{k}(C).
  \]
\end{remark}

%%% Local Variables:
%%% mode: LaTeX
%%% TeX-master: "main"
%%% End:
 % sec:comparison
\bibliographystyle{alpha}
\bibliography{biblio}

\end{document}

%%% Local Variables:
%%% mode: LaTeX
%%% TeX-master: t
%%% End:
